% =====================================================================
%  第十章  总结与体会
% =====================================================================
\chapter{总结与体会}

本次课程设计围绕 UNIX V6++ 的内存管理离散化展开，目标是在教师已有代码基础上，继续
完成页表系统、物理页框分配和进程映像管理的改造，并使系统能够在本地环境中真正运行。
最终实现将原先偏连续映像和相对映射的管理方式，改造成以 4KB 页框为单位的离散页式
管理模型——每进程独立页目录与私有用户页表承担地址映射，位示图承担页框分配，引用
计数与 COW 承担共享页的生命周期与按需复制。

通过本次课程设计，我对操作系统内存管理的理解从"概念层面的分页机制"进一步深入到
"一个可运行内核中各模块如何配合"。页表并不是孤立的数据结构，它会牵动进程切换、
异常处理、exec 参数栈、内核访问用户地址、text 共享释放以及设备 I/O 缓冲区等多条
路径；位示图也不只是一个分配器，它需要和页表项、引用计数以及释放逻辑共同维护物理
页的生命周期。这种"一个数据结构的改变会沿控制流扩散到许多看似无关的地方"的体验，
是单纯阅读教材难以获得的。

调试过程也让我认识到，内核问题往往不能只看故障发生的最后一行。一次 kernel page
fault 可能真正来自 fork 时引用计数漏加，也可能来自 text 释放范围错误，还可能来自
信号返回时恢复了错误的栈帧；一次用户程序参数丢失，也不一定是 shell 解析错误，而可能
是 exec 构造新栈时写入了错误的页表映射。只有沿着完整的执行路径分析，才能把现象和
根因正确对应起来。报告第七章列出的几个问题，正是这一认识的具象化。

此外，本次项目让我体会到工程环境对操作系统实验的重要性。中文路径、陈旧构建缓存、
Makefile 命令差异、bootloader 加载扇区数不足、磁盘镜像布局不一致，都可能造成看似
"内核崩溃"的现象。对于底层系统项目而言，稳定的构建和运行环境本身就是调试能力的一部
分——先排除工程层面的干扰，才能让真正的算法缺陷显露出来。

总体而言，本次课程设计完成了 UNIX V6++ 内存管理从连续映像思路到离散页框管理思路的
核心转变。虽然当前实现仍保留了部分历史路径，在内存不足处理、DMA 支持和自动化测试等
方面还有后续改进空间，但系统已经具备在本地构建、启动并展示页表、位示图、COW 与离散化
进程映像管理等主要机制的条件。通过这次实现与调试，我对操作系统内核中"抽象设计"与
"工程落地"之间的关系有了更具体的理解。

\clearpage
